(1) 对相位求导得 \(|\psi'_\phi\rangle=ie^{i\phi}|1\rangle/\sqrt2\)，从而 \(\langle\psi'|\psi'\rangle=1/2\)、\(|\langle\psi|\psi'\rangle|^2=1/4\)。于是单次发射单光子的 \(\boxed{F_Q=1}\)，\(\nu\) 个独立样本的局部无偏方差满足 \(\operatorname{Var}(\hat\phi)\ge1/\nu\)。该下界是局部理论，不解决相位周期性与全局先验歧义。

(2) 选择适合已知工作点 \(\phi_0\) 的投影基 \(|b_\pm\rangle=(|0\rangle\pm i e^{i\phi_0}|1\rangle)/\sqrt2\)。得到 \(p_\pm=[1\pm\sin(\phi-\phi_0)]/2\)（正负标签无关紧要），故
\[
F_{\rm cl}=\sum_\pm(\partial_\phi p_\pm)^2/p_\pm
=\cos^2(\phi-\phi_0)\bigl[1-\sin^2(\phi-\phi_0)\bigr]^{-1}=1
\]
在两结果概率非零处严格达到QFI。工作点应选择概率约为1/2以避免极端结果造成有限样本估计困难；相位未知很大时需要预估或自适应偏置。

(3) 单光子损失态为 \(\rho_\phi=\eta|\psi_\phi\rangle\langle\psi_\phi|+(1-\eta)|{\rm vac}\rangle\langle{\rm vac}|\)。两个正交区块的权重与相位无关，所以量子Fisher信息按成功概率加权：\(\boxed{F_Q(\rho_\phi)=\eta}\)。重复 \(\nu\) 次只有约 \(\eta\nu\) 次保留相位敏感样本，方差界变为 \(1/(\eta\nu)\)。如果对存活样本条件化后仍以 \(\nu\) 次全部发射作可用样本数，便违反正确的资源会计。

(4) NOON态相位的生成元在两路径分支分别为0和 \(N\)，因此其方差为 \(N^2/4\)，得 \(F_Q=N^2\)。在题设的独立损耗模型中，任意一次光子丢失即消除路径间相干，保留相位信息的分支概率是 \(\eta^N\)，故每次发射
\[
\boxed{F_Q^{\mathrm{loss}}=\eta^N N^2,\qquad F_Q^{\mathrm{loss}}/N=N\eta^N}.
\]
与 \(N\) 个独立单光子合计 \(\eta N\) 比较，单位光子增益为 \(N\eta^{N-1}\)。将 \(N\) 暂作连续量，最大化 \(N\eta^N\) 给出 \(N_*=-1/\ln\eta\)，实际需比较附近整数 \(N\ge1\)。固定 \(0<\eta<1\) 令 \(N\to\infty\) 信息趋零，不能把理想无损的 \(N^2\) 增益无条件外推到高损耗计量。
